\subsection{范阳到潼关：命令要经过粮道}
安禄山在755年起兵，见\eventref{ST-755-AN-LUSHAN-REBELLION}{安禄山起兵}。范阳的军府、东北道路和长期形成的将校关系使这次行动有了不同于普通地方骚动的基础；它仍不是一个可以由“胡将反叛”四字穷尽的原因。兵员、马匹、仓廪、朝廷任命和京师的政治信息都参与了危机。叛军向南推进时，洛阳、黄河渡口与潼关各有不同的供给与防守难题，不能把地图上的连线误写成部队必然通过的路线。

潼关失守在756年，见\eventref{LT-756-001}{潼关失守}。关隘的意义正在于它把山地、道路与粮运压缩到有限入口；一旦防线崩溃，京师不是自动失去，而是必须在仓储、人员撤离、继承名义和讯息混乱中作出新的决定。玄宗西行、肃宗在灵武立朝、回纥的援助与唐军收复两京，分别发生在不同政治与地理条件下。把它们统称为“平乱”，会抹去盟友为什么出兵、地方军为什么被倚重、两京收复后为何仍难恢复旧制。

回纥援军是唐得以重组军事选择的一部分，绝不是唐方叙事中可随意调用的背景人物。盟约、赏赐、马匹、市场和路线使双方都有自己的成本。汉文史书保存了唐朝的求援与奖赏语言，突厥—回鹘碑铭及其他材料可打开另一端，却也不自动解决日期和译名的问题。读者应把“援军”理解为两个政治共同体在危机里建立的短期合作，而非稳定从属关系。

\subsection{重建的税源从江淮出发}
战后唐廷必须养活不再主要依靠府兵的军队。刘晏整顿盐铁与转运，见\eventref{LT-758-007}{刘晏整顿盐铁转运}，不能被缩为一个聪明官员的财政策略。盐的生产、运销、商人、江淮水路、仓场与地方征收都在其中；中央要把资源送往军队和两京，也必须承认实际掌握船只、市场和地方关系的人并不都在京师。食货志、会要和通典留下政策与术语，不能证明各地都按同样办法执行。

两税法在780年施行，见\eventref{LT-780-020}{两税法施行}，通常被放进“由租庸调到两税”的整齐转换。实际的转换更像在战后户籍、土地、货币、地方强势者与军费压力之间寻找一套可征收的办法。它解决了什么，要看某个区域是否仍能登记、征纳和转输；它让谁承担代价，也不能只从法令的普遍语气推定。敦煌文书、墓志和地方材料提供的是零碎断面，正适合限制我们把制度名称当作全面现实的冲动。

\subsection{藩镇并非一种地方}
河北的成德、魏博、卢龙与江南、淮西的军政处境并不相同。它们都有军队、税源与地方中介，却处在不同的道路、市场、边防和中央关系中。平淮西的行动在819年完成，见\eventref{LT-819-037}{平淮西}；它显示唐廷仍能在特定条件下调集军力、谈判并接收一处军镇，不能据此推出所有藩镇已被同样整合。九世纪的宦官、节度使、盐利、科举官僚和地方士人，彼此争夺的是能否控制具体的人、钱粮与信息。

吐蕃进入河西及其后续变化、回鹘与唐的贸易和军事关系、南诏在西南的扩张，也让“唐后期”不能只沿长安到河北的一条线叙述。河西的敦煌材料尤其珍贵，因为它让某些军政、寺院和日常文书从中原本纪外露出；它们属于沙州的保存条件，不能代表所有边地。晚唐国家不是消失，而是在不同区域以不同的税、军、道路和谈判形式继续存在。

\subsection{黄巢之后，谁能接住城市}
874年王仙芝起事、875年黄巢转入独立行动，见\eventref{LT-874-058}{王仙芝起事}与\eventref{LT-875-059}{黄巢起事}。饥荒、盐贩、流民、地方武装和征收压力在不同地区的组合未必相同，故不能把起事者的社会构成写成一种固定身份。880年黄巢入长安，见\eventref{LT-880-062}{黄巢入长安}，使皇帝与官僚西走，城市的仓、门、市场与供给都成为新的争夺对象。

唐朝在黄巢之后得以恢复名义，却越来越依靠朱温、李克用等拥有自身军政基础的力量。907年朱温受禅并不只是唐的终点；它也是各个地区开始以不同速度处理旧唐官号、税源、军队和道路的开始。五代十国应由这些不同时钟来读，而不是由开封的朝代列表替它们盖章。
